\documentclass{article}
\usepackage{amsmath,amssymb}
\usepackage[vocab={cryptography}]{leanify}
\newtheorem{lemma}{Lemma}
\begin{document}

\section{Notation}

We write $x\bmod n$ for the residue of $x$ modulo $n$.
We write $P(a,k,n)$ for the modular power $a$ to the $k$ modulo $n$.
We write $N(p,q)$ for the RSA modulus for the primes $p$ and $q$.
We write $\varphi(p,q)$ for the RSA totient for the primes $p$ and $q$.
We write $d(p,q,e)$ for the RSA decryption exponent for the primes $p$ and $q$ and the encryption exponent $e$.
We write $\log_{g,p}(h)$ for the discrete logarithm of $h$ to the base $g$ modulo $p$.
We write $K(g,p,a,b)$ for the Diffie-Hellman key for the base $g$, the modulus $p$ and the exponents $a$ and $b$.

\section{RSA}

\begin{lemma}\label{rsa}
$N(3,11)=33$, $\varphi(3,11)=20$ and $d(3,11,3)=7$.
\end{lemma}

\begin{lemma}\label{power}
$P(2,10,1000)=24$ and $17\bmod 5=2$.
\end{lemma}

\section{Discrete logarithms}

\begin{lemma}\label{dlog}
$\log_{2,11}(8)=3$ and $K(2,11,3,4)=4$.
\end{lemma}

\end{document}
